% Einheit 2 von 5 – Parallelogramm (bank/flaechen/e2.jsonl, zusammenbau v0.1)
%% TODO Zweigzeile Teil 1: was hier gelernt wird (Satz in Schülersprache aus dem Einheitstitel) – Einheit 2
\einheitenkopf[e2]{Einheit 2 von 5 · Parallelogramm}
\zweigzeile{ab Kl. 7, je nach Buch bis Kl. 8 · keine P10-Aufgabe}
\setcounter{aufgabe}{14}

%% TODO Titel: Ich-kann-Satz fehlt in der Bank (Platzhalter: Kettenname) – Nr. 15
%% TODO Anweisung: ein Satz über der ersten Teilaufgabe fehlt in der Bank – Nr. 15
\begin{aufgabe}{Parallelogramm}
\swa{Zeichne die Höhe zur Grundseite $g$ ein und markiere den rechten Winkel.}{\parallelogramm[seiten={$g$,,,}]{4}{2.5}{60}}
%% TODO Darstellung für schwach fehlt in der Bank (flaechen-e2-k1-s1-v1; Abschnitt „Für schwache Schüler“)
\swz{Parallelogramm, $g = 9$ cm, $h = 4$ cm – Fläche?}{}
%% TODO Darstellung für schwach fehlt in der Bank (flaechen-e2-k1-s1-v2; Abschnitt „Für schwache Schüler“)
\swz{Parallelogramm, $g = 7$ m, $h = 6$ m – Fläche?}{}
%% TODO Darstellung für schwach fehlt in der Bank (flaechen-e2-k1-s1-v3; Abschnitt „Für schwache Schüler“)
\swz{Ein Fliesenstück in Parallelogrammform: $g = 15$ cm, $h = 8$ cm – Fläche?}{}
%% TODO Darstellung für schwach fehlt in der Bank (flaechen-e2-k1-s1-v4; Abschnitt „Für schwache Schüler“)
\swz{Parallelogramm, $g = 11$ dm, $h = 5$ dm – Fläche?}{}
%% TODO Darstellung für schwach fehlt in der Bank (flaechen-e2-k1-s1-v5; Abschnitt „Für schwache Schüler“)
\swz{Ein Grundstück in Parallelogrammform: $g = 25$ m, $h = 14$ m – Fläche?}{}
\swfrage{Was bleibt gleich, was ändert sich?}
%% TODO Darstellung für schwach fehlt in der Bank (flaechen-e2-k1-s2-v1; Abschnitt „Für schwache Schüler“)
\swz{Parallelogramm: Grundseite $9$ cm, schräge Seite $6$ cm, Höhe $5$ cm – Fläche?}{}
%% TODO Darstellung für schwach fehlt in der Bank (flaechen-e2-k1-s3-v1; Abschnitt „Für schwache Schüler“)
\swz{Parallelogramm, $g = 6{,}5$ cm, $h = 4$ cm – Fläche?}{}
%% TODO Darstellung für schwach fehlt in der Bank (flaechen-e2-k1-s4-v1; Abschnitt „Für schwache Schüler“)
\swz{Parallelogramm mit den Seiten $6$ cm und $4$ cm. Die Höhe zur $6$-cm-Seite ist $2$ cm, die Höhe zur $4$-cm-Seite $3$ cm. Rechne die Fläche auf beide Arten.}{}
%% TODO Darstellung für schwach fehlt in der Bank (flaechen-e2-k1-s5-v1; Abschnitt „Für schwache Schüler“)
\swz{Parallelogramm, $A = 63$ cm², $g = 9$ cm – Höhe?}{}
\swa{Ein Parallelogramm hat die Grundseite $g = 6{,}4$ cm und die Höhe $h = 3{,}5$ cm. Zeichne in der Skizze die Höhe zu $g$ ein, berechne die Fläche und begründe die Formel: Wie wird aus dem Parallelogramm ein Rechteck?}{\parallelogramm[seiten={$g$,,,}]{4.5}{2.5}{60}}
\end{aufgabe}

%% TODO Titel: Ich-kann-Satz fehlt in der Bank (Platzhalter: Kettenname) – Nr. 16
%% TODO Anweisung: ein Satz über der ersten Teilaufgabe fehlt in der Bank – Nr. 16
\begin{aufgabe}{Umfang aus den Seiten}
%% TODO Darstellung für schwach fehlt in der Bank (flaechen-e2-k2-s1-v1; Abschnitt „Für schwache Schüler“)
\swz{Parallelogramm mit den Seiten $9$ cm und $6$ cm – Umfang?}{}
\end{aufgabe}

%% TODO Titel: Ich-kann-Satz fehlt in der Bank (Platzhalter: Kettenname) – Nr. 17
%% TODO Anweisung: ein Satz über der ersten Teilaufgabe fehlt in der Bank – Nr. 17
\begin{aufgabe}{Parallelogramm – Fehler finden, Begründen, Anwendung}
\swz[3]{Parallelogramm: Grundseite $7$ cm, schräge Seite $5$ cm, Höhe $4$ cm. Nils rechnet: \rechnung{A &= 7 \cdot 5 \\ A &= 35 \text{ cm}^2} Wo ist der Fehler?}{}
\swfrage{Warum ist die Fläche eines Parallelogramms Grundseite mal Höhe? Begründe mit Abschneiden und Anlegen.}
\swz[3]{Ein Parkstreifen hat die Form eines Parallelogramms: $45$ m lang (Grundseite) und $5{,}5$ m breit (Höhe). Er wird gepflastert, $1$ m² kostet $38$ €. Was kostet das Pflaster?}{}
\end{aufgabe}

\uebersichtskasten{Fläche: Grundseite mal Höhe – die Höhe steht senkrecht auf der Grundseite, nicht die schräge Seite. \\ g = 8 cm, h = 5 cm: A = 8 · 5 = 40 cm² \\ Warum: Dreieck links abschneiden, rechts anlegen – es entsteht ein Rechteck.}
